\fichatitulo{26 · REVISÃO}{Preprint v1 · 2023}{Panorama de RAG}
{\small Gao et al. \textit{Retrieval-Augmented Generation for Large Language Models: A Survey}. Preprint v1 · 2023. \href{https://arxiv.org/pdf/2312.10997v1}{Fonte consultada}.\par}

\campo{Problema e objetivo}
Como organizar estratégias de geração aumentada por recuperação?

\campo{Principal contribuição · síntese interpretativa}
Sistematiza os paradigmas Naive, Advanced e Modular RAG e as etapas de recuperação, geração e ampliação.

\campo{Método / natureza da evidência}
Revisão e taxonomia de técnicas, incluindo processamento anterior e posterior à recuperação.

\campo{Resultado ou argumento principal}
Ajuda a localizar alterações de pipeline e relações entre seus componentes, sem fornecer uma configuração universalmente ótima.

\campo{Excertos-chave · seleção literal}
\begin{itemize}
\excerto{1}{Naive RAG, Advanced RAG, and Modular RAG}{Resumo.}
\excerto{2}{pre-retrieval and post-retrieval}{Advanced RAG.}
\end{itemize}

\campo{Limitações e análise crítica}
Análise da ficha: o trecho sobre ARES nesta v1 contém descrição e comparação que precisam ser confrontadas com o artigo primário. Não usar a revisão como evidência da superioridade de um avaliador.

\campo{Aproveitamento na dissertação · proposta}
Fundamentação: adotar a taxonomia; para resultados e métricas, recorrer a ARES, RAGAS e demais fontes primárias.

\registro{Fonte consultada nesta preparação: Resumo, arquitetura, Advanced RAG e trecho de avaliação; consulta parcial. Chave: \texttt{gao2023ragSurvey}.}
